\scanpage{73}
% archive foliation: 72; typescript printed page number follows
\begin{center}
- 18 -
\end{center}
\[
(2)\qquad f_i(x_1,\ldots,x_p\,;\,v_1(x_1,\ldots,x_p),\ldots,v_q(x_1,\ldots,x_p))=0
\quad(1\leq i\leq q)
\]
\uline{dans un voisinage de} $(a_1,\ldots,a_p)$, \uline{on a}
$u_i=v_i$ $(1\leq i\leq q)$ \uline{dans un voisinage de}
$(a_1,\ldots,a_p)$.

Par translation, on peut se ramener au cas où $a_j=0$ pour $1\leq j\leq p$
et $b_i=0$ pour $1\leq i\leq q$,. Posons pour $1\leq i\leq q$
\[
(3)\quad
f_i(x_1,\ldots,x_p;y_1,\ldots,y_q)
=f_{i0}+\sum_{j=1}^q f_{ij}y_j
+\sum_{(n_j)}f_{i;n_1\ldots n_q}y_1^{n_1}\ldots y_q^{n_q}
\]
les $f_{i0}$, $f_{ij}$ et $f_{i;n_1\ldots n_q}$ étant des séries entières par
rapport aux $x_k$ $(1\leq k\leq p)$ sommables dans un polycylindre ouvert $P$
de centre l'origine dans $K^p$, $f_{i0}$ n'ayant pas de terme constant ;
en outre $f_i$ est sommable dans $P\times Q$, où $Q$ est un polycylindre de centre
l'origine dans $K^q$, et dans la seconde sommation des seconds membres de (3)
on a toujours $\sum_{j=1}^q n_j\geq2$.
Par hypothèse, la matrice $\underline F=(f_{ij})$ est holomorphe dans $P$ et
inversible à l'origine donc elle est inversible dans un polycylindre ouvert
contenu dans $P$, et son inverse $\underline G=(g_{ij})$ est holomorphe dans ce
polycylindre ; nous supposerons donc chacune des $g_{ij}$ développées en série
entière sommable dans un polycylindre ouvert $P_1\subset P$ de centre l'origine ;
si on pose alors $g_i=\sum_{j=1}^q g_{ij}f_j$, on aura
\[
(4)\quad
g_i(x_1,\ldots,x_p;y_1,\ldots,y_q)
=-h_{i0}+y_i-\sum_{(n_j)}h_{i;n_1\ldots n_q}y_1^{n_1}\ldots y_q^{n_q}
\]
où $h_{i0}$ et les $h_{i;n_1\ldots n_q}$ sont des séries entières par rapport
aux $x_k$, sommables dans $P_1$, $h_{i0}$ n'ayant pas de terme constant, et la
somme étant étendue aux $(n_j)$ tels que $\sum_{j=1}^q n_j\geq2$.
La première partie du théorème revient à prouver qu'il existe $q$ séries entières
sans terme constant $u_i$ $(1\leq i\leq q)$ par rapport aux $x_k$, sommables dans
un voisinage de l'origine, et telles que l'on ait
\[
(5)\qquad
u_i=h_{i0}+\sum_{(n_j)}h_{i;n_1\ldots n_q}u_1^{n_1}\ldots u_q^{n_q}
\qquad\text{pour }1\leq i\leq q.
\]
